% ============ 第2章 §2.8 内函数与外函数(原书 pp.144-147) ============
\section{内函数与外函数}\label{sec:2.8}
\index{内函数!(单位圆板上的)}
\index{外函数!(单位圆板上的)}
\index{内函数!(单位圆板上的)奇异内函数}
\index{外函数!(单位圆板上的)奇异内函数}
\index{因式分解定理}


定理~\ref{thm:2.6.4} 给出了属于 $H^{p}$ 的函数 $f$ 的分解 $f = Bg$,其中 $g$ 在 $D$ 内无零点,本节进一步分解 $g$.事实上,利用对某个实测度 $\mu \in M[-\pi, \pi]$ 的 Herglotz 积分表示(见 2.2 节)
\begin{equation*}
\log g(z) = \frac{1}{2\pi}\int_{-\pi}^{\pi}\frac{e^{it} + z}{e^{it} - z}\dd\mu(t) \quad (z \in D).
\end{equation*}
此时,对 $\mu$ 作 Lebesgue 分解 $\mu = \mu_{a} + \mu_{s}$,$\mu_{s}$ 是正测度.这里 $\mu_{a}$ 关于 Lebesgue 测度绝对连续且 $\mu_{s}$ 奇异.定理~\ref{thm:2.8.3} 表明 $\mu$ 的这个分解给出 $g$ 的分解.下面的定理~\ref{thm:2.8.1} 是命题~\ref{prop:2.2.1} 的推广,但对证明因数分解定理~\ref{thm:2.8.3} 很重要.

\begin{theorem}\label{thm:2.8.1}
设 $0 < p \leq \infty$.若 $f \in H^{p}$,则
\begin{equation*}
\log|f(re^{i\theta})| \leq \frac{1}{2\pi}\int_{-\pi}^{\pi}P(r, \theta - t)\log|f(e^{it})|\dd t.
\end{equation*}
\end{theorem}

\begin{proof}
由定理~\ref{thm:2.6.4} 可分解 $f = Bg$,故若对 $g$ 证明,则由 $|B(z)| < 1$($z \in D$)且 $|B(e^{i\theta})| = 1$ a.e. $\theta$ 可对 $f$ 证明不等式.因此可以认为 $f$ 在 $D$ 内无零点.此时 $\log|f(z)|$ 在 $D$ 上调和,故当 $r < \rho < 1$ 时,
\begin{equation*}
\log|f(\rho re^{i\theta})| = \frac{1}{2\pi}\int_{-\pi}^{\pi}P(r, \theta - t)\log|f(\rho e^{it})|\dd t.
\end{equation*}
由系~\ref{cor:2.7.4},
\begin{equation*}
\lim_{\rho \to 1}\int_{-\pi}^{\pi}P(r, \theta - t)\log^{+}|f(\rho e^{it})|\dd t = \int_{-\pi}^{\pi}P(r, \theta - t)\log^{+}|f(e^{it})|\dd t,
\end{equation*}
且由 Lebesgue 积分的 Fatou 引理,
\begin{equation*}
\lim_{\rho \to 1}\int_{-\pi}^{\pi}P(r, \theta - t)\log^{-}|f(\rho e^{it})|\dd t \geq \int_{-\pi}^{\pi}P(r, \theta - t)\log^{-}|f(e^{it})|\dd t,
\end{equation*}
由此可导出本定理.
\end{proof}

$Q \in H^{\infty}$ 称为内函数,是指 $|Q(e^{i\theta})| = 1$ a.e. $\theta$;$F \in H^{p}$ 称为外函数,是指
\begin{equation*}
\log|F(re^{i\theta})| = \frac{1}{2\pi}\int_{-\pi}^{\pi}P(r, \theta - t)\log|F(e^{it})|\dd t.
\end{equation*}
定理~\ref{thm:2.8.1} 表明不等式总是成立.若 $F$ 是外函数,则 $\log|F| \in L^{1}$,故 $F$ 在 $D$ 内无零点.由定理~\ref{thm:2.6.1},Blaschke 积 $B$ 是内函数.下面的定理~\ref{thm:2.8.2} 与在 $D$ 内无零点的内函数和外函数的 Herglotz 积分表示有关.设 $\mu \in M[-\pi, \pi]$ 为关于 Lebesgue 测度奇异的正测度,且
\begin{equation*}
S(z) = \exp\left(-\int_{-\pi}^{\pi}\frac{e^{it} + z}{e^{it} - z}\dd\mu(t)\right),
\end{equation*}
则 $\mu'(\theta) = 0$ a.e. $\theta$,且
\begin{equation*}
\log|S(z)| = -\int_{-\pi}^{\pi}P(r, \theta - t)\dd\mu(t),
\end{equation*}
故由定理~\ref{thm:2.3.2},$|S(e^{i\theta})| = 1$ a.e. $\theta$.因此 $S$ 是内函数,且 $S$ 在 $D$ 内无零点.$S$ 称为奇异内函数.

\begin{theorem}\label{thm:2.8.2}
(1) 若 $Q$ 是在 $D$ 内无零点的内函数,则存在奇异的有界正测度 $\mu$,满足 $\mu(\{-\pi\}) = \mu(\{\pi\})$,使得
\begin{equation*}
Q(z) = e^{i\gamma}\exp\left\{-\int_{-\pi}^{\pi}\frac{e^{it} + z}{e^{it} - z}\dd\mu(t)\right\}
\end{equation*}
可以表示.这里 $\gamma$ 是实常数.

(2) 若 $F$ 是外函数,则
\begin{equation*}
F(z) = e^{i\gamma}\exp\left\{\int_{-\pi}^{\pi}\frac{e^{it} + z}{e^{it} - z}\log|F(e^{it})|\dd t\right\}
\end{equation*}
可以表示.这里 $\gamma$ 是实常数.
\end{theorem}

\begin{proof}
(1) 由最大值原理与 Fatou 定理~\ref{thm:2.3.2},$|Q(z)| < 1$($z \in D$),故 $-\log|Q(z)|$ 是 $D$ 上的正调和函数.由系~\ref{cor:2.4.2},存在正测度 $\mu \in M[-\pi, \pi]$ 使得
\begin{equation*}
-\log|Q(z)| = \frac{1}{2\pi}\int_{-\pi}^{\pi}P(r, \theta - t)\dd\mu(t) \quad (z \in D).
\end{equation*}
由 $|Q(e^{i\theta})| = 1$ a.e. $\theta$ 与定理~\ref{thm:2.3.2},$\mu'(\theta) = 0$ a.e. $\theta$.因此 $\mu$ 是奇异正测度.

令
\begin{equation*}
q(z) = \frac{1}{2\pi}\int_{-\pi}^{\pi}\frac{e^{it} + z}{e^{it} - z}\dd\mu(t),
\end{equation*}
则 $\operatorname{Re}(q(z)) = -\log|Q(z)|$($z \in D$).因此存在实常数 $\gamma$ 使得 $-\log Q(z) = q(z) + i\gamma$.

(2) 当 $F$ 是外函数时,令
\begin{equation*}
G(z) = \frac{1}{2\pi}\int_{-\pi}^{\pi}\frac{e^{it} + z}{e^{it} - z}\log|F(e^{it})|\dd t,
\end{equation*}
则由外函数的定义,$\operatorname{Re}G(z) = \log|F(z)|$($z \in D$).因此存在实常数 $\gamma$ 使得 $\log F(z) = G(z) + i\gamma$.
\end{proof}

\begin{theorem}\label{thm:2.8.3}
设 $0 < p \leq \infty$.若 $f \in H^{p}$ 不恒等于零,则 $f$ 可以唯一分解为 $f = BSF$.这里 $B$ 是 Blaschke 积,$S$ 是奇异内函数,$F$ 是 $H^{p}$ 中的外函数.
\end{theorem}

\begin{proof}
由定理~\ref{thm:2.6.4},可分解 $f = Bg$,$g$ 在 $D$ 内无零点.由定理~\ref{thm:2.5.3},$f \in L^{p}$ 且 $\log|f| \in L^{1}$.

令
\begin{equation*}
F(z) = \exp\left\{\frac{1}{2\pi}\int_{-\pi}^{\pi}\frac{e^{it} + z}{e^{it} - z}\log|f(e^{it})|\dd t\right\},
\end{equation*}
则 $|g(e^{i\theta})| = |F(e^{i\theta})|$ a.e. $\theta$.因此令 $g(z)/F(z) = Q(z)$,$Q$ 是在 $D$ 内无零点的内函数,由定理~\ref{thm:2.8.2} 的 (1),$Q(z) = e^{i\gamma}S(z)$.把 $F$ 换成 $e^{i\gamma}F$,定理得证.
\end{proof}
